% Original ebook customization on the unmodified ElegantBook v4.7 class.
% Template: https://github.com/ElegantLaTeX/ElegantBook (LPPL 1.3c+).
\usepackage{mathtools,tabularx,longtable,array,ragged2e}
\usepackage{pgfplots}
\pgfplotsset{compat=1.18}
\usetikzlibrary{arrows.meta,positioning,calc,shapes.geometric,fit,decorations.pathreplacing}
\usepackage{bookmark}
\usepackage[nameinlink,noabbrev]{cleveref}
\geometry{paperwidth=190mm,paperheight=250mm,left=20mm,right=20mm,
 top=21mm,bottom=22mm,headheight=16pt,headsep=7mm,footskip=11mm}
\newif\ifbookwindowsfonts
\bookwindowsfontsfalse
\IfFontExistsTF{SimSun}{
 \IfFontExistsTF{Microsoft YaHei}{
  \IfFontExistsTF{KaiTi}{\bookwindowsfontstrue}{}
 }{}
}{}
\ifbookwindowsfonts
 \setCJKmainfont[BoldFont=Microsoft YaHei,ItalicFont=KaiTi]{SimSun}
 \setCJKsansfont{Microsoft YaHei}
 \setCJKmonofont{Microsoft YaHei}
\else
 \setCJKmainfont[BoldFont=FandolSong-Bold,ItalicFont=FandolKai-Regular]{FandolSong-Regular}
 \setCJKsansfont[BoldFont=FandolHei-Bold]{FandolHei-Regular}
 \setCJKmonofont{FandolFang-Regular}
\fi
\definecolor{ink}{HTML}{203C44}
\definecolor{teal}{HTML}{276A73}
\definecolor{gold}{HTML}{A87934}
\definecolor{paper}{HTML}{FAF8F2}
\colorlet{structurecolor}{ink}
\colorlet{main}{teal}
\colorlet{second}{gold}
\colorlet{third}{teal}
\hypersetup{pdftitle={扩散模型：从分布直觉到多模态生成},
 pdfauthor={扩散模型学习手册},pdfsubject={DDPM、score、SDE、概率流、Flow Matching 与多模态生成},
 pdfkeywords={Diffusion Models,DDPM,Flow Matching,DiT,中文教材},
 linkcolor=teal,citecolor=teal,urlcolor=teal,bookmarksopenlevel=1,
 pdfpagelayout=OneColumn,pdfstartview=FitH}
\setcounter{tocdepth}{1}
\setcounter{secnumdepth}{2}
\setstretch{1.32}
\setlength{\parskip}{3pt plus 1pt}
\setlength{\emergencystretch}{2em}
\raggedbottom
\allowdisplaybreaks[2]
\AtBeginDocument{\setlength{\abovedisplayskip}{9pt plus 2pt minus 2pt}
 \setlength{\belowdisplayskip}{9pt plus 2pt minus 2pt}}
\setlist[itemize]{leftmargin=2em,itemsep=3pt,topsep=5pt}
\setlist[enumerate]{leftmargin=2em,itemsep=3pt,topsep=5pt}
\renewcommand{\arraystretch}{1.2}
\newcolumntype{Y}{>{\RaggedRight\arraybackslash}X}
\fancyhf{}
\fancyhead[L]{\small\color{ink}扩散模型}
\fancyhead[R]{\small\color{ink}\nouppercase{\rightmark}}
\fancyfoot[C]{\small\color{ink}\thepage}
\renewcommand{\headrulewidth}{0.4pt}
\fancypagestyle{plain}{\fancyhf{}\fancyfoot[C]{\small\color{ink}\thepage}
 \renewcommand{\headrulewidth}{0pt}\renewcommand{\headrule}{}}
\tcbset{bookbox/.style={enhanced,breakable,sharp corners,
 boxrule=0pt,borderline west={1.4pt}{0pt}{teal},colback=teal!4,
 colframe=teal,coltitle=ink,colbacktitle=teal!7,
 fonttitle=\sffamily\bfseries,fontupper=\normalfont,
 left=10pt,right=10pt,top=7pt,bottom=7pt,before skip=10pt,after skip=10pt}}
\newtcolorbox{intuitionbox}[1]{bookbox,title={直觉\quad #1}}
\newtcolorbox{keybox}[1]{bookbox,title={要点\quad #1}}
\newtcolorbox{cautionbox}[1]{bookbox,borderline west={1.4pt}{0pt}{gold},colback=gold!5,title={辨析\quad #1}}
\newtcolorbox{examplebox}[2][]{bookbox,title={例子\quad #2},#1}
\newtcolorbox{derivationbox}[1]{bookbox,title={推导\quad #1}}
\newtcolorbox{exercisebox}[1]{bookbox,colback=black!2,title={练习\quad #1}}
\newcommand{\readingguide}[1]{\begin{keybox}{本章阅读路线}#1\end{keybox}}
\newcommand{\chaptertakeaway}[1]{\begin{keybox}{带走的理解}#1\end{keybox}}
\newcommand{\term}[2]{\hyperref[#2]{\textcolor{teal}{#1}}}
\newcommand{\E}{\mathbb{E}}
\newcommand{\R}{\mathbb{R}}
\newcommand{\N}{\mathcal{N}}
\DeclareMathOperator{\KL}{KL}
\newcommand{\dd}{\mathop{}\!\mathrm{d}}
\newcommand{\norm}[1]{\left\lVert#1\right\rVert}
\newcommand{\ip}[2]{\left\langle#1,#2\right\rangle}
\newcommand{\pdata}{p_{\mathrm{data}}}
\newcommand{\pnoise}{p_{\mathrm{noise}}}
\newcommand{\score}{s}
\newcommand{\eps}{\epsilon}
\newcommand{\vtheta}{v_\theta}
\newcommand{\utheta}{u_\theta}
\newcommand{\grad}{\nabla}
\DeclareMathOperator{\divg}{div}
\DeclareMathOperator{\tr}{tr}
\DeclareMathOperator{\Var}{Var}
\DeclareMathOperator{\Cov}{Cov}
\DeclareMathOperator{\SNR}{SNR}
\DeclareMathOperator{\argmin}{arg\,min}
\crefname{chapter}{章}{章}\crefname{section}{节}{节}
\crefname{equation}{式}{式}\crefname{figure}{图}{图}\crefname{table}{表}{表}
\ExecuteBibliographyOptions{sorting=none,maxbibnames=6,maxcitenames=2,doi=false,isbn=false,eprint=false}
\setlength{\bibitemsep}{7pt}
\AtEveryBibitem{\clearfield{urlday}\clearfield{urlmonth}\clearfield{urlyear}}
